\subsubsection{Depuração e reparo}

A síntese parcial de depuração cobre 59 estudos de depuração e reparo de programas, com 126 comparações extraídas nessa tarefa. No
sentido da métrica, 82 favorecem a condição com a técnica, 22 a desfavorecem e 3 empatam. As outras 19 não têm
direção definida, por não terem sentido declarado ou os dois valores. Em 46 estudos, ao menos uma técnica focal ou
componente tem avaliação do pacote, e em 38, ao menos uma tem ingrediente isolado ou contraste pareado. Outros 8 estudos não têm
comparação numérica nessa tarefa. Os agrupamentos abaixo são desta revisão, e um estudo pode entrar em mais de um.
As contagens por agrupamento vêm das sínteses parciais. Os 3 estudos de depuração incluídos na última leva da
etapa 3 estão no fim da subseção.

\paragraph{Feedback de testes e execução}
Treze estudos devolvem ao modelo, em nova rodada, resultados de testes, de depurador, de verificadores ou contraexemplos. Só quatro isolam o feedback como ingrediente (S62, S114, S412, S613). S65 traz só rodadas cumulativas, e ao menos cinco estudos o avaliam apenas dentro de pacotes.
% AFIRMACAO: A6
No ARHE, as observações reais do depurador deram ao AutoSD patch plausível em 73\% das execuções, contra 63\% com observações previstas pelo próprio LLM (S62-C2).
% AFIRMACAO: A7
Com ChatGPT no Defects4J, o retorno das falhas de testes elevou de 36 para 62 os bugs corrigidos (S412-C3); no DS-1000, retirar o conhecimento de bug derivado dos testes reduziu o ANF de 104 para 83 (S613-C3).

Há também nulos e inversões.
% AFIRMACAO: A8
Com InCoder 1B, as estruturas de reparo P1 e P2 empataram em 0,0\% (S114-C2); em S2164, a taxa AC subiu de 72\% para 73,53\% no LeetCode, mas, segundo a heterogeneidade registrada na ficha, caiu de 31\% para 29\% no Codeforces (S2164-C2).
O número de chamadas raramente é igualado (o ThinkRepair admite até 25 sessões; S412-C1), e em S65 as localizações fornecidas vêm do gabarito, isto é, são informação oráculo.

\paragraph{Informação diagnóstica no prompt}
Vinte estudos inserem uma única vez mensagem do compilador, tipo de defeito, testes que falham, traços, relato de bug ou localização.
% AFIRMACAO: A9
Em S647, a localização perfeita, informação oráculo, elevou as correções corretas de 34 para 44 (S647-C1); trocar os exemplos por teste disparador ou por relato de bug levou-as a 59 e 72, mudança que retira e acrescenta ingredientes ao mesmo tempo (S647-C2; S647-C3).
% AFIRMACAO: A10
Com Qwen2.5-Coder:7B em 40 programas do QuixBugs, descrição, linha defeituosa ou saída dos testes elevaram a taxa de reparo de 62,5\% para 70\%, enquanto o PDG e o arquivo de testes a reduziram para 55\% e 57,5\% (S1674-C1; S1674-C2; S1674-C5; S1674-C3; S1674-C4).
A linha defeituosa, nesse estudo, vem da versão correta (oráculo).
% AFIRMACAO: A11
Em S1678, acrescentar stack trace ao prompt de localização reduziu os bugs localizados de 47 para 32, e stack trace com asserção, para 41, em contagem cuja unidade de agregação não é informada (S1678-C1; S1678-C2).
% AFIRMACAO: A12
Em S233, as mensagens de execução aumentaram a Pass Rate nos bugs de sintaxe e de referência e a reduziram nos de lógica ($-$4,3 com GPT-4 e com GPT-3.5) (S233-C1); em S1008, mais contexto não teve efeito na compilação ($\chi^2 = 0{,}21$; $p = 0{,}649$) (S1008-C1).
% AFIRMACAO: A13
Entre os acréscimos favoráveis estão os traços de execução, de 154 para 177 correções corretas (S1304-C2), e as ferramentas de busca, de 14 para 23 patches plausíveis em 100 bugs (S303-C2).
Uma leitura possível, sem isolamento, é que o efeito dependa do tipo de informação e da categoria do bug, já que informação estrutural ou volumosa aparece entre as pioras de S1674 e S1678.

\paragraph{Contexto recuperado}
Quinze estudos inserem código chamador ou chamado, conhecimento de API, design rationales, memória de reparos ou padrões recuperados; oito têm contraste de ingrediente (S275, S435, S613, S783, S979, S2046, S2119, S1943).
% AFIRMACAO: A14
Com GPT-3.5 no RepoBugs, o RLCE (variante one-shot) elevou a taxa de reparo correto de 0,2258, com contexto só da função defeituosa, para 0,5968 (S275-C1).
% AFIRMACAO: A15
Na ablação de S2046, com GPT-4o no Defects4J V1.2 e localização perfeita, a busca interna por dependências elevou as correções corretas de 85 para 116, e a busca externa de padrões, para 108 (S2046-C2; S2046-C3).
% AFIRMACAO: A16
Mais contexto, porém, não equivale a ganho: injetar de uma vez todas as camadas de conhecimento deu Pass@1 de 0,38, contra 0,5 da injeção progressiva (S821-C2), e o DeepK-S ficou em 0,7449, contra 0,7466 do ThinkRepair (S2119-C2).
Os orçamentos diferem (1 candidata no BaseChatGPT e até 300 no RelRepair; S979-C1), e vários estudos supõem localização perfeita (S315, S821, S979, S2046).

\paragraph{Exemplos}
Oito estudos comparam few-shot fixo ou seleção de pares defeito-correção.
% AFIRMACAO: A17
Há melhoras com seleção de exemplos por diagnóstico (de 0,76 para 0,82; S25-C2), com dois exemplos contra completamento zero-shot, em pacote (de 6,7\% para 20,3\%; S53-C1), e com memória episódica (de 44,7\% para 48,3\%; S783-C3).
% AFIRMACAO: A18
Três comparações de exemplos fixos não são favoráveis, com queda de 0,783 para 0,643 com DeepSeek Coder 1.3B (S1598-C1), de 205 para 200 plausíveis com GPT-4 (S2057-C1) e empate em 11 (S1565-C4); em S552, a implementação de referência e o programa mais próximo reduziram os programas corrigidos de 606 para 576 e 578 (S552-C3; S552-C4).

\paragraph{Raciocínio explícito}
No reparo, as sínteses parciais registram três comparações de ingrediente, em três estudos, e seis de pacote, em quatro estudos.
% AFIRMACAO: A19
Duas das três comparações de ingrediente pioram: o CoT teve 0,5161, contra 0,5645 da instrução Detail (S275-C2), e pedir explicação antes do refinamento deu pass@1 de 40,37\%, contra 40,64\% (S354-C2); a melhora de 36 para 57 em S412 inclui seleção de exemplos com raciocínio, e não só a instrução de raciocinar (S412-C2).
% AFIRMACAO: A20
Em pacotes, com Qwen3 Coder Next em circuitos quânticos, o Pass@1 caiu de 95\% (Structured) para 45\% com CoT e 63\% com ReAct (S1429-C1; S1429-C2), e com DeepSeek Coder 1.3B o CoT reduziu o overall\_score de 0,783 para 0,613 (S1598-C2).
Essas pioras vêm de um modelo pequeno e de um domínio específico, e em S1429 o conteúdo dos prompts não é descrito.

\paragraph{Decomposição, planejamento e agentes}
Onze estudos avaliam pacotes com análise de causas, planos, ferramentas ou vários agentes; há ablações em S810, S654, S303 e S1304.
% AFIRMACAO: A21
No T3, a taxa de reparo foi de 46,7\%, contra 37,6\% sem a etapa de plano e 37,3\% sem a análise de causas, com modelo da ablação não informado (S810-C2; S810-C3).
% AFIRMACAO: A22
Em S303, retirar a máquina de estados reduziu os patches plausíveis de 23 para 18 (S303-C3), e em S1304 retirar o debate multiagente reduziu as correções corretas de 177 para 158 (S1304-C3).
% AFIRMACAO: A23
Entre sistemas, o TraceRepair com DeepSeek-V3.2 corrigiu 392 bugs, contra 291 do REINFIX com GPT-4o, mas, sob o mesmo GPT-3.5, corrigiu 224 contra 241 segundo a heterogeneidade registrada (S1304-C1); o RepairAgent superou o ChatRepair por 164 a 162, relação que se inverte no subtotal Defects4J v1.2 (S303-C1).
Com número de chamadas e backbones distintos, os ganhos de pacote não se atribuem a um mecanismo único.

\paragraph{Instruções e formato}
Oito estudos variam templates, restrições, persona, táticas de influência ou modo de edição.
% AFIRMACAO: A24
Em S2054, o prompt guiado por taxonomia elevou a CER de 50,4\% para 57,3\% com GPT-4o-mini (S2054-C4); em S1565, um estilo alternativo de comentário reduziu as falhas corrigidas de 10 para 5 (S1565-C3); e em S1556 as táticas de influência não tiveram efeito principal no SWE-bench Verified (S1556-C2).

\paragraph{Estudos extraídos depois da síntese parcial}
% AFIRMACAO: A113
Em S1869, com GPT-4 Turbo no TFix, acrescentar ao prompt o identificador da regra violada elevou o CodeBLEU de
0,444 para 0,496, e combinações de dois tipos de informação chegaram a 0,513 e 0,528 (S1869-C1 a S1869-C3).
% AFIRMACAO: A114
Em S1715, um prompt que induz memorização levou o GPT-3.5 a gerar, no Defects4J, patches idênticos ao de referência
em 78,76\% dos patches corretos, resultado que o estudo usa como indício de memorização do \textit{benchmark}
(S1715-C1). S1624 só tem resultados descritivos.

\paragraph{Lacunas}
\begin{itemize}
  \item O feedback de execução só é isolado como ingrediente em 4 dos 13 estudos, e o efeito do número de rodadas não tem suporte numérico (S1429).
  \item O raciocínio explícito reúne poucas comparações, com sinais opostos e domínios restritos.
  \item Linguagens de hardware (S342, S393, S1132), circuitos quânticos (S1429), OCaml (S2130), acessibilidade web (S2135) e reparo de segurança (S2129, com 7 casos; S391-C2) dependem de um estudo ou de conjuntos sem repetição.
  \item Das 126 comparações da tarefa, 19 não têm direção definida; S1620 imprime só variação relativa, sem acurácias absolutas.
  \item Comparadores reaproveitados de outros artigos (S321-C1, S783-C1, S2046), orçamentos desiguais e métricas de plausibilidade limitam a comparação.
  \item A contaminação não é informada em 27 dos 59 estudos da síntese parcial, e nenhuma comparação mede produtividade, custo ou segurança em uso real.
\end{itemize}
